\documentclass[11pt,a4paper]{article}
\usepackage[utf8]{inputenc}
\usepackage{amsmath,amssymb,amsthm}
\usepackage{listings}
\usepackage{geometry}
\usepackage{hyperref}
\geometry{margin=2.5cm}

\lstset{
  language=C++,
  basicstyle=\ttfamily\small,
  keywordstyle=\bfseries,
  breaklines=true,
  frame=single,
  numbers=left,
  numberstyle=\tiny,
  tabsize=4,
  showstringspaces=false
}

\title{IOI 2007 -- Miners}
\author{Solution}
\date{}

\begin{document}
\maketitle

\section{Problem Statement}
$N$ food carts arrive sequentially, each carrying one of 3 food types (M, F, G). There are 2 mines. Each cart must be assigned to one mine. When a mine receives a cart, the \emph{bonus} is the number of distinct food types among its last (up to) 3 deliveries. Maximize the total bonus.

\section{Solution}

\subsection{DP Formulation}
The bonus at a mine depends only on its \textbf{last two deliveries} (plus the new one). Encode the state as $(a_1, a_2, b_1, b_2)$ where $a_1, a_2$ are the last two deliveries to mine~1 and $b_1, b_2$ to mine~2. Each coordinate takes values in $\{0,1,2,3\}$ (three food types plus a ``none'' sentinel for fewer than 2 past deliveries). This gives $4^4 = 256$ states.

\textbf{Transition.} For cart $i$ with food type $c$:
\begin{itemize}
    \item \emph{Send to mine 1:} New state $(c, a_1, b_1, b_2)$. Bonus $= |\{c, a_1, a_2\} \setminus \{3\}|$.
    \item \emph{Send to mine 2:} New state $(a_1, a_2, c, b_1)$. Bonus $= |\{c, b_1, b_2\} \setminus \{3\}|$.
\end{itemize}

\subsection{Complexity}
\begin{itemize}
    \item \textbf{Time:} $O(256 \cdot 2 \cdot N) = O(N)$.
    \item \textbf{Space:} $O(256) = O(1)$ (rolling two DP layers).
\end{itemize}

\section{C++ Solution}

\begin{lstlisting}
#include <bits/stdc++.h>
using namespace std;

int main(){
    ios::sync_with_stdio(false);
    cin.tie(nullptr);

    int N;
    cin >> N;
    string s;
    cin >> s;

    auto encode = [](char c) -> int {
        if(c == 'M') return 0;
        if(c == 'F') return 1;
        return 2;
    };

    // Pack state (a1, a2, b1, b2) into a single int
    auto pack = [](int a1, int a2, int b1, int b2) -> int {
        return a1 * 64 + a2 * 16 + b1 * 4 + b2;
    };

    // Count distinct food types among c, x, y (ignoring sentinel 3)
    auto bonus = [](int c, int x, int y) -> int {
        bool seen[3] = {};
        seen[c] = true;
        if(x < 3) seen[x] = true;
        if(y < 3) seen[y] = true;
        return seen[0] + seen[1] + seen[2];
    };

    vector<int> dp(256, -1);
    dp[pack(3, 3, 3, 3)] = 0;

    for(int i = 0; i < N; i++){
        int c = encode(s[i]);
        vector<int> ndp(256, -1);
        for(int st = 0; st < 256; st++){
            if(dp[st] == -1) continue;
            int a1 = (st / 64) % 4, a2 = (st / 16) % 4;
            int b1 = (st / 4) % 4,  b2 = st % 4;

            // Send to mine 1
            int ns1 = pack(c, a1, b1, b2);
            int val1 = dp[st] + bonus(c, a1, a2);
            ndp[ns1] = max(ndp[ns1], val1);

            // Send to mine 2
            int ns2 = pack(a1, a2, c, b1);
            int val2 = dp[st] + bonus(c, b1, b2);
            ndp[ns2] = max(ndp[ns2], val2);
        }
        dp = ndp;
    }

    cout << *max_element(dp.begin(), dp.end()) << "\n";
    return 0;
}
\end{lstlisting}

\section{Notes}
The problem is a clean DP exercise. The state space is $4^4=256$ because each mine needs only its two most recent deliveries (plus the sentinel for ``no delivery yet''). The bonus function counts distinct types in a window of size~3 (current delivery plus last~2). Since the state space is constant, the algorithm runs in $O(N)$ time.

\end{document}
